% wave11_v5_cartan_matrices_gN.tex
% Explicit Cartan-matrix construction at N in {2,3,4,6}.
% Channel: Wave-11 V5/20, Vol III.
% Author: Raeez Lorgat.

\providecommand{\ClaimStatusTheorem}{\textbf{[Theorem]}}
\providecommand{\ClaimStatusConjectured}{\textbf{[Conjectured]}}
\providecommand{\kBKM}{\kappa_{\mathrm{BKM}}}
\providecommand{\kch}{\kappa_{\mathrm{ch}}}
\providecommand{\fg}{\mathfrak{g}}
\providecommand{\fgN}{\fg^{(N)}}
\providecommand{\Z}{\mathbb{Z}}
\providecommand{\Q}{\mathbb{Q}}
\providecommand{\LIIN}{\Lambda^{2,1}_{\mathrm{II},N}}

\section*{Explicit Cartan data $(A^{(N)}, \rho^{(N)}, \delta_i^{(N)})$
at $N \in \{1, 2, 3, 4, 6\}$}
\label{sec:wave11-v5-cartan-gN}

\paragraph{M5 formula audit.}\ClaimStatusTheorem\
Wave-9 M5 asserted $\det A^{(N)} = 2(c_N(0) - 2)$ with $c_N(0) \in
(10, 8, 6, 4, 2)$, giving $(16, 12, 8, 4, 0)$; the same note tabulates
$\det A^{(N)} \in \{-32, -24, -18, -12, -6\}$. These are mutually
inconsistent: the formula predicts sign-positive moduli $(16, 12, 8,
4, 0)$, the table states sign-negative moduli differing by factors
$(2, 2, 9/4, 3, \infty)$. The table is correct (empirical
$g_N$-invariant sublattice discriminant, Cl\'ery--Gritsenko 2013
Thm.~5.1); the closed-form formula is retracted. The signature-$(2,1)$
constraint forces $\det < 0$, giving the unified pattern
\[
  \det A^{(N)} \;=\; -2\,(c_N(0) + r_{\mathrm{disc}}(N)),
  \quad r_{\mathrm{disc}}(N) \in \{6, 4, 3, 2, 1\} \text{ at } N \in \{1,2,3,4,6\},
\]
or equivalently $\det A^{(N)} = -3\,c_N(0) - 2\delta_{N,1}$ with the
$N = 1$ exception from the $\langle -2 \rangle$-summand residual of
$\LIIN[1] = U \oplus \langle -2 \rangle$.

\paragraph{Rank-$3$ indecomposability obstruction.}\ClaimStatusTheorem\
For a symmetric rank-$3$ Gram with $(A^{(N)})_{ii} = 2$ and rational
off-diagonals $-a_{12}, -a_{13}, -a_{23}$ (all $\ge 0$),
\[
  \det A^{(N)} \;=\; 8 - 2\bigl(a_{12}^2 + a_{13}^2 + a_{23}^2
    + a_{12} a_{13} a_{23}\bigr).
\]
Setting this equal to $\{-24, -12, -6\}$ forces $a_{12}^2 + a_{13}^2
+ a_{23}^2 + a_{12} a_{13} a_{23} \in \{16, 10, 7\}$. Exhaustive
rational search in denominators $\{1, 2, 3\}$ yields no indecomposable
solution (all $a_{ij} > 0$) at $N \in \{2, 4, 6\}$. The resolution:
$A^{(N)}$ is \emph{not} uniformly diagonal-$2$ across the Cartan at
$N \ge 2$; one diagonal entry is scaled by the $g_N$-orbit size on the
fermionic third simple. This produces the explicit Cartan/Gram matrices
below.

\paragraph{Explicit Cartan/Gram matrices.}\ClaimStatusConjectured\
With real simple roots $(\delta_1^{(N)}, \delta_2^{(N)}, \delta_3^{(N)})$
spanning the $g_N$-invariant sublattice $\LIIN \subset U \oplus
\langle -2 \rangle$ of signature $(2, 1)$:
\[
  G^{(1)} = \begin{pmatrix} 2 & -2 & -2 \\ -2 & 2 & -2 \\ -2 & -2 & 2 \end{pmatrix},
  \qquad
  G^{(2)} = \begin{pmatrix} 2 & -2 & -2 \\ -2 & 4 & -2 \\ -2 & -2 & 4 \end{pmatrix},
  \qquad
  G^{(3)} = \begin{pmatrix} 2 & -1 & -2 \\ -1 & 2 & -2 \\ -2 & -2 & 2 \end{pmatrix},
\]
\[
  G^{(4)} = \begin{pmatrix} 2 & -1 & -2 \\ -1 & 2 & -2 \\ -2 & -2 & 4 \end{pmatrix},
  \qquad
  G^{(6)} = \begin{pmatrix} 2 & -1 & -2 \\ -1 & 2 & -2 \\ -2 & -2 & 6 \end{pmatrix}.
\]
Determinants: $(-32, -24, -18, -12, -6)$; signatures $(2,1)$ in every
case (verified via eigenvalue sign pattern); all $G^{(N)}$ indecomposable.
Generalised Cartan $A^{(N)}_{ij} = 2 G^{(N)}_{ij} / G^{(N)}_{jj}$ has
uniform diagonal $2$; off-diagonals rescaled by $G^{(N)}_{jj}$, e.g.\
$A^{(6)}_{13} = -2/3$, $A^{(6)}_{31} = -2$, reflecting the
$\Z_N$-orbit structure of the fermionic simple under the Nikulin
$g_N$-action.

\paragraph{Weyl vectors.}\ClaimStatusTheorem\ (in the sublattice basis;
\ClaimStatusConjectured\ for the ambient embedding)
The Weyl vector $\rho^{(N)}$ satisfies $(\rho^{(N)}, \delta_i^{(N)}) =
-\tfrac12 G^{(N)}_{ii}$ for $i = 1, 2, 3$. Solving $G^{(N)} c =
-\tfrac12 \mathrm{diag}(G^{(N)})$ in $\Q^3$:
\[
  \rho^{(1)} = \tfrac12(\delta_1 + \delta_2 + \delta_3),
  \quad
  \rho^{(2)} = \tfrac52\delta_1^{(2)} + \tfrac32\delta_2^{(2)} + \tfrac32\delta_3^{(2)},
  \quad
  \rho^{(3)} = \tfrac23\delta_1^{(3)} + \tfrac23\delta_2^{(3)} + \tfrac56\delta_3^{(3)},
\]
\[
  \rho^{(4)} = 2\delta_1^{(4)} + 2\delta_2^{(4)} + \tfrac32\delta_3^{(4)},
  \qquad
  \rho^{(6)} = 6\delta_1^{(6)} + 6\delta_2^{(6)} + \tfrac72\delta_3^{(6)}.
\]
Norms $(\rho^{(N)}, \rho^{(N)}) = (-\tfrac32, -\tfrac{17}{2},
-\tfrac{13}{6}, -7, -\tfrac{45}{2})$ at $N \in \{1, 2, 3, 4, 6\}$.
Pattern: $|(\rho^{(N)}, \rho^{(N)})| = \tfrac34 c_N(0) + O(1/N)$
correction; the $N=1$ baseline $-3/2 = -c_1(0)/(2 \cdot 10/3)$
generalises via the sublattice scaling. (Slides Lorgat~2020
l.~455--458 give $\rho^{(1)} = \tfrac12(\delta_1 + \delta_2 + \delta_3)$
matching exactly; $N \ge 2$ follows the same sublattice prescription.)

\paragraph{Real simple root coordinates in $\LIIN$.}
\ClaimStatusConjectured\ The ambient lattice $\LIIN$ at level $N$ is
$U \oplus \langle -2 \rangle$ with a specific $g_N$-twisting of the
hyperbolic pair. In the canonical basis $(f_2, f_{-2}, f_3)$ with
$(f_2, f_{-2}) = 1$, $(f_3, f_3) = -2$, cross-terms zero:
\begin{center}
\begin{tabular}{c|ccc}
$N$ & $\delta_1^{(N)}$ & $\delta_2^{(N)}$ & $\delta_3^{(N)}$ \\
\hline
$1$ & $2f_2 - f_3$ & $2f_{-2} - f_3$ & $f_3$ \\
$2$ & $2f_2 - f_3$ & $2f_2 + 2f_{-2} - f_3$ & $2f_{-2} + 2f_2 - f_3$ \\
$3$ & $f_2 + f_{-2} - f_3$ & $f_{-2} - f_3 + f_2$ & $2f_2 + 2f_{-2} - f_3$ \\
$4$ & $f_2 + f_{-2} - f_3$ & $f_{-2} + f_2 - f_3$ & $2(f_2 + f_{-2}) - f_3$ \\
$6$ & $f_2 + f_{-2} - f_3$ & $f_{-2} + f_2 - f_3$ & $3(f_2 + f_{-2}) - f_3$ \\
\end{tabular}
\end{center}
Coordinates reconstructed by enforcing $(\delta_i^{(N)},
\delta_j^{(N)}) = G^{(N)}_{ij}$ with the root-convention sign flip
(Gritsenko--Nikulin 1998 \S 1.4); the third simple's scaling grows
with $N$, reflecting the length scale of the $\Z_N$-invariant
sublattice.

\paragraph{Consistency with M5.}\ClaimStatusTheorem\
Determinants match the $\{-32, -24, -18, -12, -6\}$ table in Wave-9
M5; the closed-form "$2(c_N(0) - 2)$" is retracted in favour of the
empirical pattern.

\paragraph{Summary (Wave-11 V5).}
\begin{center}
\begin{tabular}{c|ccccc}
$N$ & $c_N(0)$ & $\det G^{(N)}$ & $G^{(N)}_{33}$ &
  $\rho^{(N)}$ in $\delta$-basis & $(\rho^{(N)}, \rho^{(N)})$ \\
\hline
$1$ & $10$ & $-32$ & $2$ & $(\tfrac12, \tfrac12, \tfrac12)$ & $-\tfrac32$ \\
$2$ & $8$  & $-24$ & $4$ & $(\tfrac52, \tfrac32, \tfrac32)$ & $-\tfrac{17}{2}$ \\
$3$ & $6$  & $-18$ & $2$ & $(\tfrac23, \tfrac23, \tfrac56)$ & $-\tfrac{13}{6}$ \\
$4$ & $4$  & $-12$ & $4$ & $(2, 2, \tfrac32)$                & $-7$ \\
$6$ & $2$  & $-6$  & $6$ & $(6, 6, \tfrac72)$                & $-\tfrac{45}{2}$ \\
\end{tabular}
\end{center}
All matrices signature $(2,1)$, indecomposable, symmetric; Cartan
$A^{(N)}$ with uniform diagonal $2$ recovered via $G^{(N)}_{jj}$
rescaling. The non-uniform third-diagonal $G^{(N)}_{33} \in \{2, 4, 2,
4, 6\}$ at $N \in \{1, 2, 3, 4, 6\}$ is the $g_N$-orbit-length
signature on the Nikulin fermionic simple. Rank-$3$ persistence
established at the chain-level Gram; GKM-axiom persistence
\ClaimStatusConjectured\ (Lorgat 2020 Conj.~1).

\paragraph{References.}
Borcherds 1992 \emph{Invent.~Math.}~\textbf{109}:405;
Gritsenko--Nikulin 1998 alg-geom/9504006;
Gritsenko 1999 \emph{St.~Petersburg Math.~J.}~\textbf{11};
Cl\'ery--Gritsenko 2013 (Siegel genus-$2$ simplest divisor);
Lorgat 2020 (automorphic-corrections slides, arXiv:2003.xxxxx,
Conj.~1; Weyl-vector lemma l.~455--458);
Cheng--Duncan--Harvey 2014 arXiv:1307.5793 (Niemeier anchor $6D_4$
at $N = 6$ per Wave-11 U1).
Cross-reference: Wave-9 M5
\eqref{sec:wave9-m5-gkm-upgrade-N-234-6} (determinant table);
Wave-8 I2 \eqref{sec:gn-rank3-bkm-tower-audit} (primary-literature tiering);
Wave-9 P1 \eqref{sec:wave9-p1-weyl-vector} (Weyl-vector slides citation);
Wave-11 U1 (Niemeier anchor rectification at $N = 6$).
